\label{chap:83-07}

La conexión se construye a partir de la clasificación orbital completa y de
su primera elevación coinductiva. El desarrollo siguiente fija las tres
historias iniciales, las fronteras orientadas y la cadena
\(w_6\to w_{30}\to R_{36}\) antes de introducir la holonomía nonádica.

\begin{HMTScientificOwnerSubordinated}
\input{ampliaciones_sucesoras_20260824/parte_i_ii/owners/U008_orbitas_elevacion_r36}
\end{HMTScientificOwnerSubordinated}

% Unidad U009. Conexión nonádica conjunta; redacción autónoma.

\section{Conexión nonádica conjunta, compresión y memoria terminal}
\label{p12-83-c7-s1:chap:conexion-nonadica-conjunta}

Las nueve fases que aparecen después de \(R_{36}\) no son nueve pruebas
independientes. Son las nueve componentes locales de una única conexión
sobre estados enriquecidos. Su composición alrededor de una vuelta define
una holonomía. El retorno de la etiqueta de fase conserva la continuidad de
prefijo, acarreo, frontera, sombra de Witt y memoria.

\subsection{Transporte del estado enriquecido por la holonomía nonádica y actualización de memoria}

En el nivel \(K\), escribimos
\[
 \widehat x_K=
 (B_K,C_K,p_K,c_K,\Sigma_K,W_K,g(K),q_{{\rm mem},K}),
\]
donde:
\begin{enumerate}
 \item \(B_K=(u_K^\pi,u_K^e,u_K^\varphi)\) es el bloque visible de tres
       palabras de seis trits;
 \item \(C_K\) es el cilindro ternario compatible;
 \item \(p_K\) es el prefijo completo;
 \item \(c_K\) es el acarreo acumulado;
 \item \(\Sigma_K\) es la firma conjunta de frontera;
 \item \(W_K\) es la sombra de incidencia de Witt;
 \item \(g(K)=1+((K-1)\bmod9)\) es la etiqueta pública de fase;
 \item \(q_{{\rm mem},K}\in\mathbb Z_{\geq0}\) es la memoria vertical no
       acotada.
\end{enumerate}
La coordenada dodecafásica es derivada:
\[
 \beta_K=[q_{{\rm mem},K}]_{12}.
\]
No se sustituye \(q_{\rm mem}\) por \(\beta\), pues la reducción módulo
doce perdería el número de vueltas.

\subsection{Relaciones locales y holonomía}

Para cada \(K\), sea
\[
 \mathcal R_K\subseteq
 \widehat X_K\times\widehat X_{K+1}
\]
la relación de extensión compatible entre estados completos. La composición
relacional de una vuelta es
\begin{equation}
 \Gamma_{9,K}
 :=\mathcal R_{K+8}\circ\cdots\circ\mathcal R_{K+1}\circ\mathcal R_K
 :\widehat X_K\longrightarrow\widehat X_{K+9}.
\label{p12-83-c7-s1:eq:gamma9-composicion-local}
\end{equation}

\begin{definition}[Holonomía nonádica]
\label{p12-83-c7-s1:def:holonomia-nonadica}
La holonomía de una vuelta es
\[
 \Gamma_9:=\operatorname{Hol}_{C_9}(\nabla_{\rm TPK}),
\]
cuya realización en el nivel \(K\) es
\(\Gamma_{9,K}\). Las nueve relaciones \(\mathcal R_K,\ldots,
\mathcal R_{K+8}\) son cartas de esta única conexión.
\end{definition}

\begin{theorem}[Retorno de fase sin reinicio]
\label{p12-83-c7-s1:thm:retorno-nonadico-sin-reinicio}
Si \((\widehat x_K,\widehat x_{K+9})\in\Gamma_{9,K}\), entonces
\[
 g(K+9)=g(K),
 \qquad
 q_{{\rm mem},K+9}=q_{{\rm mem},K}+1,
\]
y
\[
 \beta_{K+9}=\beta_K+1\pmod{12}.
\]
En general, \(\widehat x_{K+9}\neq\widehat x_K\).
\end{theorem}

\begin{proof}
La primera identidad procede de la definición de \(g\). Una vuelta completa
incrementa por definición el contador de vueltas no acotado. Reducir esa
identidad módulo doce da la tercera. Las restantes coordenadas se
transportan por la composición \cref{p12-83-c7-s1:eq:gamma9-composicion-local}; ninguna
ley del TPK las reinicializa. Por tanto, la igualdad de fase no implica
igualdad del estado completo.
\end{proof}

\subsection{Clasificación exhaustiva de las \(9\) fases del Topological Phase Kernel}

Los bloques visibles y los censos locales de la primera vuelta son
\[
\begin{array}{c|rrrrr|ccc}
K&N_\pi&N_e&N_\varphi&N_{\rm loc}&N_{\rm ext}
 &u^\pi&u^e&u^\varphi\\\hline
1&378&422&349&3&2&012222&121221&112202\\
2&238&368&388&2&5&010211&102011&121020\\
3&184&284&173&5&4&002111&012222&010210\\
4&207&207&122&4&1&110221&102011&010200\\
5&39&151&151&1&1&222220&021222&102011\\
6&110&110&69&1&3&111201&201202&221021\\
7&81&29&81&3&3&212121&222210&200221\\
8&59&59&59&3&2&200121&212212&110110\\
9&43&43&43&2&1&100100&020112&010122
\end{array}.
\]

\begin{theorem}[Funciones de las nueve componentes]
\label{p12-83-c7-s1:thm:funciones-nueve-componentes}
En la vuelta anterior, las nueve componentes desempeñan las funciones
siguientes:
\begin{enumerate}
 \item apertura local de tres orientaciones y primera resolución;
 \item separación binaria con máximo local de cinco extensiones;
 \item máxima multiplicidad local de cinco estados;
 \item contracción de cuatro estados locales a una extensión;
 \item primera saturación fuerte, con una sola continuación;
 \item presente local único y apertura triple del futuro;
 \item fibra especular triple;
 \item sincronización de los tres censos en \(59\);
 \item cierre orientado: dos estados locales de entrada y una relación que,
       tras el horizonte adicional, conserva una prolongación.
\end{enumerate}
\end{theorem}

\begin{proof}
Cada afirmación se lee de los censos \(N_{\rm loc},N_{\rm ext}\), de la
igualdad o desigualdad de \(N_\pi,N_e,N_\varphi\) y de la relación de
extensión explícita en las fases novena, décima y undécima que se demostrará
en la unidad siguiente. La tabla enumera toda la vuelta; no se ha escogido
un subconjunto de fases.
\end{proof}

La quinta componente coincide con la primera frontera conjunta
\[
 R_{36}=(222220|021222|102011)^T.
\]
La novena componente contiene
\[
 B_9=(100100|020112|010122).
\]
La presencia de dos estados locales en la fase novena y de tres extensiones
desde el estado orientado pertenece a capas diferentes; no son censos
contradictorios.

La holonomía actúa sobre el estado enriquecido completo. Antes de publicar
su cociente modal se construyen, por tanto, la involución que intercambia
las fibras de \(\pi\) y \(e\), el eje fijo de \(\varphi\), los invariantes
ternarios y la resolución local \(3\to1\) determinada por el segundo
horizonte.

\begin{HMTScientificOwnerSubordinated}
\input{ampliaciones_sucesoras_20260824/parte_i_ii/owners/U010_espejo_seleccion_tres_a_uno}
\end{HMTScientificOwnerSubordinated}

\subsection{Transferencia conjunta de la fibra de 468 emisiones}

La factorización del catálogo es
\[
 \mathcal U_6^{\rm cat}\cong
 \mathfrak S_+\times\mathfrak S_\times,
 \qquad
 |\mathfrak S_+|=18,\qquad |\mathfrak S_\times|=26.
\]
Para una función \(f\) sobre el catálogo, sean
\[
 (E_+f)(s,e)=\frac1{18}\sum_{s'\in\mathfrak S_+}f(s',e),
\]
y
\[
 (E_\times f)(s,e)=\frac1{26}
 \sum_{e'\in\mathfrak S_\times}f(s,e').
\]
Los proyectores \(E_+\) y \(E_\times\) conmutan.

Sea
\[
 \tau=(+1,-1,0,+1,-1,0,+1,-1,0),
 \qquad
 P_{+1}=E_+,quad P_{-1}=E_\times,quad P_0=I.
\]
En \(\mathcal U_6^{\rm cat}\times C_9\) definimos la transferencia de fase
por
\begin{equation}
 \mathcal K_{\rm ph}\bigl((u,r),(v,[r+1]_9)\bigr)
 :=P_{\tau(r)}(u,v),
 \label{p12-83-c7-s1:eq:u009-kph-explicita}
\end{equation}
y la hacemos cero para los demás pares de fases. Así, la permutación cíclica
y el operador aplicado en cada uno de los nueve pasos quedan determinados.

\begin{proposition}[Renovación conjunta tras una vuelta]
\label{p12-83-c7-s1:prop:transferencia-conjunta-nueve}
\[
 \mathcal K_{\rm ph}^{\,9}=E_+E_\times.
\]
En particular, sobre una masa puntual,
\[
 (E_+E_\times)(u,v)=\frac1{18\cdot26}=\frac1{468}.
\]
\end{proposition}

\begin{proof}
El producto ordenado sobre una vuelta es
\[
 P_{\tau(8)}P_{\tau(7)}\cdots P_{\tau(0)}
 =I E_\times E_+ I E_\times E_+ I E_\times E_+
 =E_+^3E_\times^3I^3=E_+E_\times.
\]
La segunda igualdad usa la conmutatividad y la última la idempotencia.
El segundo resultado es la distribución uniforme sobre el producto de las
dos firmas.
\end{proof}

Esta identidad expresa una renovación conjunta, no nueve filtraciones
independientes.

\subsection{Descomposición de la holonomía \(U_{\Gamma_9}\) en compresión visible \(T\) y componente ortogonal \(\eta\): \(I-T^*T=\eta^*\eta\)}

Sean \(\mathcal H_{\rm vis}\) y \(\mathcal H_{\rm full}\) espacios de
Hilbert para la realización visible y completa. Sea
\[
 J:\mathcal H_{\rm vis}\longrightarrow\mathcal H_{\rm full}
\]
una inclusión isométrica y
\(U_{\Gamma_9}\) una realización isométrica de la holonomía. Definimos
\[
 T:=J^*U_{\Gamma_9}J,
 \qquad
 \eta:=(I-JJ^*)U_{\Gamma_9}J.
\]

\begin{theorem}[Descomposición compresión--expresión]
\label{p12-83-c7-s1:thm:conexion-compresion-expresion}
Se cumplen
\[
 \boxed{U_{\Gamma_9}J=JT+\eta,}
\]
\[
 J^*\eta=0,
 \qquad
 \boxed{I-T^*T=\eta^*\eta.}
\]
\end{theorem}

\begin{proof}
Descomponemos \(U_{\Gamma_9}J\) mediante las proyecciones ortogonales
\(JJ^*\) e \(I-JJ^*\):
\[
 U_{\Gamma_9}J
 =JJ^*U_{\Gamma_9}J+(I-JJ^*)U_{\Gamma_9}J
 =JT+\eta.
\]
La ortogonalidad da \(J^*\eta=0\). Como \(J\) y
\(U_{\Gamma_9}\) son isometrías,
\[
 I=J^*U_{\Gamma_9}^*U_{\Gamma_9}J
 =T^*T+\eta^*\eta.
\]
Reordenar prueba la última identidad.
\end{proof}

\(T\) es la compresión visible; \(\eta\) es la expresión complementaria.
Ninguna de las dos debe llamarse contracción estricta. Para esa función se
reserva un operador distinto.

\subsection{Contracción estricta y telescopía terminal}

Sea \(V_9\) una isometría y sea \(q\) un parámetro con \(0<q<1\). Definimos
\[
 M_9=qV_9,
 \qquad
 D_9=(I-M_9^*M_9)^{1/2}.
\]
Entonces \(\|M_9\|=q<1\) y
\[
 I-M_9^*M_9=D_9^*D_9.
\]

\begin{theorem}[Identidad de Stein con término terminal]
\label{p12-83-c7-s1:thm:stein-telescopia-terminal}
Para todo \(L\geq1\),
\[
 \boxed{
 I=
 \sum_{r=0}^{L-1}M_9^{*r}D_9^*D_9M_9^r
 +M_9^{*L}M_9^L.}
\]
Más generalmente, si \(S_0\) satisface
\(S_0-M_9^*S_0M_9=D_9^*D_9\), entonces
\[
 \boxed{
 S_0=
 \sum_{r=0}^{L-1}M_9^{*r}D_9^*D_9M_9^r
 +M_9^{*L}S_0M_9^L.}
\]
\end{theorem}

\begin{proof}
Multiplicamos la identidad de defecto a izquierda por \(M_9^{*r}\) y a
derecha por \(M_9^r\). Al sumar para \(r=0,\ldots,L-1\), los términos
interiores telescopan y permanecen el inicial y el terminal. La misma prueba
se aplica a la ecuación general de Stein.
\end{proof}

El último término no se borra en profundidad finita. Su desaparición en un
límite requiere una prueba de convergencia; no es parte de la identidad
algebraica.

\subsection{Cociclo de acarreo y potencias relacionales}

Para flechas composables \(\gamma_1,\gamma_2\), el acarreo satisface
\[
 \operatorname{Carr}(\gamma_2\gamma_1)
 =\operatorname{Carr}(\gamma_2)H(\gamma_1)
 +Y(\gamma_2)\operatorname{Carr}(\gamma_1).
\]
La fórmula conserva la contribución de la primera etapa después de
transportarla por la segunda.

Las vueltas de mayor longitud se escriben con índices desplazados:
\begin{align*}
 \Gamma_{27,K}
 &=\Gamma_{9,K+18}\circ\Gamma_{9,K+9}\circ\Gamma_{9,K},\\
 \Gamma_{54,K}
 &=\Gamma_{9,K+45}\circ\cdots\circ\Gamma_{9,K},\\
 \Gamma_{108,K}
 &=\Gamma_{9,K+99}\circ\cdots\circ\Gamma_{9,K}.
\end{align*}
Sólo después de declarar un operador global sobre la unión graduada puede
abreviarse esta notación como \(\Gamma_9^3,\Gamma_9^6,\Gamma_9^{12}\).

\subsection{Estado nonádico completo: fase, acarreo, frontera, ruta y memoria tras la prolongación}

La conexión de una vuelta se registra por la quíntupla
\[
 \mathfrak N_9=
 \bigl(\Gamma_9,T,\eta,\operatorname{Carr},S_{\rm term}\bigr),
\]
donde \(S_{\rm term}=M_9^{*L}S_0M_9^L\) en una profundidad finita
\(L\). Un funtor de aplicación que utilice la conexión debe declarar su
compatibilidad con los cinco campos. Conmutar sólo con la fase visible no
constituye naturalidad respecto del estado nonádico completo.

\subsection{Obstrucciones de la holonomía \(\Gamma_9\): retorno \(9\to1\), índices de composición y conservación de fase, acarreo, frontera, ruta y memoria}

\begin{criterion}[Criterios de refutación de la conexión]
La construcción queda refutada si:
\begin{enumerate}
 \item se presentan las nueve fases como filtros independientes;
 \item el retorno \(9\to1\) reinicia \(q_{\rm mem}\), prefijo o carry;
 \item se llama contracción estricta a \(T\);
 \item falla alguna identidad de la
       \cref{p12-83-c7-s1:thm:conexion-compresion-expresion};
 \item se elimina el término terminal a profundidad finita;
 \item \(\Gamma_{27,K}\) se compone sin desplazar los índices;
 \item un funtor de aplicación ignora alguno de los cinco campos de
       \(\mathfrak N_9\).
\end{enumerate}
\end{criterion}

Los desarrollos completos que siguen reconstruyen primero el emisor finito,
su factorización y sus microfibras; después despliegan la involución
especular, las tres extensiones de \(B_9\), la unicidad a dos horizontes y la
inversa Bockstein--Hensel que conserva la memoria durante el retorno de fase.

\begin{HMTScientificOwnerSubordinated}
\input{sections/hmt/03b_alfabeto_factorizacion}
\end{HMTScientificOwnerSubordinated}

\begin{HMTScientificOwnerSubordinated}
\input{sections/hmt/14_numero_heisenberg_y_d108}
\end{HMTScientificOwnerSubordinated}
